\documentclass[aspectratio=169,11pt]{beamer}
\usetheme{metropolis}
\usepackage{minted,booktabs}
\setminted{fontsize=\scriptsize,breaklines,bgcolor=black!3}
\setbeamertemplate{navigation symbols}{}
\title{From a shell to an HPC result}
\author{SURF course materials}
\date{}
\begin{document}
\maketitle
\begin{frame}{What you will be able to do}
\begin{enumerate}
\item Find your way around the system and select software.
\item Submit a job with resources that match the program.
\item Check the answer before interpreting its performance.
\item Place files deliberately and express workflow dependencies.
\end{enumerate}
\end{frame}
\begin{frame}{One path through the exercises}
\begin{center}
Shell $\longrightarrow$ modules $\longrightarrow$ Slurm $\longrightarrow$ results\par\bigskip
Results $\longrightarrow$ parallelism $\longrightarrow$ I/O $\longrightarrow$ workflows
\end{center}
\medskip
Keep the first job small enough to understand. The backup slides hold full
scripts and details; the module decks provide the extended reference.
\end{frame}
\begin{frame}{Start with where you are}
Identify your login node, home directory and course directory.
\begin{itemize}
\item The login node is for preparation and light interactive work.
\item A batch job runs on resources selected by the scheduler.
\item Home, shared scratch and job-local temporary storage serve different needs.
\end{itemize}
Explain which machine executes each command.
\end{frame}
\begin{frame}{Select software before submitting work}
\begin{itemize}
\item Modules select compilers, libraries and applications.
\item Record the modules that the job requires inside its script.
\item On the SRC teaching cluster, stack aliases select the EESSI environment.
\item Check the site's instructions; an alias does not reproduce its hardware.
\end{itemize}
\end{frame}
\begin{frame}[fragile]{First result: run Python through Slurm}
From \texttt{course-module-fundamentals-of-hpc/hands-on}:
\begin{minted}{bash}
sbatch python.sh
squeue -u "$USER"
cat slurm-JOBID.out
\end{minted}
Replace \texttt{JOBID} with the number returned by \texttt{sbatch}.
Expect the compute hostname and \texttt{sum: 499999500000}.
Find the actual runnable script in backup.
\end{frame}
\begin{frame}{Resources describe how the program runs}
\begin{description}
\item[Tasks] Separate processes, commonly MPI ranks.
\item[CPUs per task] Cores available to each process, including OpenMP threads.
\item[Nodes] Machines across which tasks are placed.
\end{description}
One sequential process does not use extra CPUs simply because they were requested.
\end{frame}
\begin{frame}[fragile]{Threads: check pi, then compare runtimes}
From the same hands-on directory:
\begin{minted}{bash}
bash compilepi
sbatch --cpus-per-task=2 pi.sh
\end{minted}
The script runs with one and two threads inside the allocation.
Both answers should be approximately \texttt{3.141592654}.
Discuss correctness, elapsed time and overhead separately.
\end{frame}
\begin{frame}{MPI: split data and combine the answer}
\begin{itemize}
\item Use the hybrid-jobs AXPY and dot-product examples.
\item Try a size not divisible by the number of ranks.
\item AXPY prints \texttt{y[0] = 202}; the global dot product must match the serial result.
\item Explain which process owns each element and where results are combined.
\end{itemize}
Use the SRC notes for small teaching allocations.
\end{frame}
\begin{frame}{Files belong where their lifetime requires}
\begin{description}
\item[Home] Keep scripts and small persistent material.
\item[Shared scratch] Share active data across compute nodes, subject to site policy.
\item[Job temporary directory] Use \texttt{\$TMPDIR} inside the compute job.
\end{description}
The filesystem exercise writes bounded data and reports actual bytes.
Inspect free space before increasing the input size.
\end{frame}
\begin{frame}{Many tasks need an explicit dependency}
For the CSV exercise, compute each dataset's mean, then aggregate.
\begin{itemize}
\item Aggregation waits for every required producer to succeed.
\item It consumes exactly the selected input set.
\item A failed producer must make the workflow fail.
\end{itemize}
Compare the QCG teaching variants before using the workflow alternatives.
\end{frame}
\begin{frame}{Compare engines through observable results}
QCG-PilotJob, HyperQueue, Snakemake and FireWorks express dependencies differently.
\begin{itemize}
\item Use the same datasets and check the same CSV means.
\item Repeat with fewer inputs to expose retained output.
\item Inspect task states and the command's exit status.
\item Compare performance only with the worker count and allocation stated.
\end{itemize}
\end{frame}
\begin{frame}{Explain your result to a neighbour}
\begin{enumerate}
\item Which machine ran the program, with which software?
\item Which tasks and threads did it use?
\item What expected value checked the answer?
\item Which files and dependencies did it need?
\end{enumerate}
\end{frame}
\appendix
\section{Backup: runnable scripts and reference}
\begin{frame}[fragile]{Python job: included from the runnable file}
\inputminted{bash}{../course-module-fundamentals-of-hpc/hands-on/python.sh}
\end{frame}
\begin{frame}[fragile]{Compile pi: included from the runnable file}
\inputminted{bash}{../course-module-fundamentals-of-hpc/hands-on/compilepi}
\end{frame}
\begin{frame}[fragile]{OpenMP job: included from the runnable file}
\inputminted{bash}{../course-module-fundamentals-of-hpc/hands-on/pi.sh}
\end{frame}
\begin{frame}[fragile]{The pi calculation}
\inputminted[fontsize=\tiny]{c}{../course-module-fundamentals-of-hpc/hands-on/pi.c}
\end{frame}
\begin{frame}{Independent numerical checks}
\begin{itemize}
\item AXPY: $y_0=202$ for the exercise's initial values.
\item Dot product: $25\sin(N)\sin(N-1)/\sin(1)$ for the given initialization.
\item Matrix multiplication: all entries are $2N$, so the checksum is $2N^3$.
\item Binary filesystem output: $8N$ bytes for $N$ doubles on the teaching platform.
\end{itemize}
Check uneven MPI partitions as well as sizes divisible by the rank count.
\end{frame}
\begin{frame}{SRC resource examples}
\begin{tabular}{lll}
\toprule
Exercise & Processes & CPUs per process\\\midrule
Sequential AXPY, dot, matrix & 1 & 1\\
MPI AXPY & 2 & 1\\
Hybrid dot (three nodes) & 6 & 1\\
OpenMP pi & 1 & 2\\\bottomrule
\end{tabular}
\medskip
These are correctness allocations. Read the module's SRC notes for complete
submission overrides and ensure the nodes have enough memory.
\end{frame}
\begin{frame}{Workflow exercise versions are teaching steps}
The QCG versions expose races, ordering and dependency choices. A finishing
launcher alone does not establish a correct result.
\begin{itemize}
\item Explain the dependency graph in each version.
\item Inspect the final version's aggregate against the selected inputs.
\item Run alternatives in separate directories to preserve their outputs.
\end{itemize}
\end{frame}
\begin{frame}{Extended module references}
\begin{itemize}
\item Fundamentals: site architecture, accounting and Slurm directives.
\item Hybrid jobs: serial, MPI and OpenMP source and submission scripts.
\item Filesystems: navigation, permissions and I/O exercises.
\item Performance analysis: profiling tools and interpretation.
\item QCG-PilotJob and workflow alternatives: orchestration exercises.
\end{itemize}
The concise teaching path selects a first lesson; the extended decks remain
available for site-specific questions and deeper classes.
\end{frame}
\end{document}
